\section*{Bab \ref{ch:b2:fragment-orbitals} --- Orbital Fragmen dan Molekul Poliatomik}

\begin{solution}{exo:b2:fragment-orbitals:1}
$h_1 + h_2$ dan $h_1 - h_2$. Bidang $xz$, yang membagi dua sudut itu dan
tegak lurus molekul, menukar kedua hidrogen: $h_1 - h_2$ berganti tanda,
$h_1 + h_2$ tidak.
\end{solution}

\begin{solution}{exo:b2:fragment-orbitals:2}
Enam orbital valensi (2s dan tiga 2p oksigen, dua 1s hidrogen), delapan
elektron valensi, empat orbital terisi.
\end{solution}

\begin{solution}{exo:b2:fragment-orbitals:3}
Tingkat-tingkat valensi yang terisi adalah satu orbital $a_1$ dan sehimpunan
tiga orbital $t_2$ yang terdegenerasi: dua energi, dua pita. Pita $t_2$
mengambil elektron dari tiga orbital yang memuat enam elektron,
dibandingkan dengan satu orbital dan dua elektron untuk $a_1$.
\end{solution}

\begin{solution}{exo:b2:fragment-orbitals:4}
Ikatan $\pi$ mensyaratkan kedua orbital p karbon sejajar, yang menempatkan
kedua bidang \ce{CH2} pada satu bidang yang sama.
\end{solution}

\begin{solution}{exo:b2:fragment-orbitals:5}
Amonia kehilangan sebuah elektron dari pasangan elektron bebasnya, orbital
yang sebagian besar nonikatan pada nitrogen; metana dari \omterm{def:b2:diatomic-mos:bonding}{orbital ikatan}
$t_2$. Elektron nonikatan terletak lebih tinggi, lebih dekat ke tingkat
atom, daripada elektron ikatan: elektron itu lebih mudah diambil.
\end{solution}

\begin{solution}{exo:b2:fragment-orbitals:6}
Air linear mempunyai dua orbital p nonikatan terisi yang terdegenerasi;
pembengkokan menurunkan salah satunya dan sedikit menaikkan sebuah tingkat
ikatan: keuntungan bersih dengan delapan elektron. \ce{BeH2} hanya
mempunyai empat, yang mengisi kedua \omterm{def:b2:diatomic-mos:bonding}{orbital ikatan}; orbital yang akan
diturunkan oleh pembengkokan kosong, sementara tingkat-tingkat ikatan
kehilangan tumpang-tindih: bentuk linear paling baik.
\end{solution}

\begin{solution}{exo:b2:fragment-orbitals:7}
Setiap ikatan C--H atau C--C pada karbon di sebelah pusat kationik dapat
sejajar dengan orbital p kosong dan memberikan stabilisasi sekitar
$2\beta^2/\Delta$. \ce{CH3+} tidak mempunyai tetangga seperti itu; kation
etil, isopropil, dan tert-butil mempunyai satu, dua, dan tiga gugus metil
di sebelah pusat itu, masing-masing menawarkan sebuah ikatan yang sejajar:
kestabilannya bertambah menurut urutan itu.
\end{solution}

\begin{solution}{exo:b2:fragment-orbitals:8}
Untuk dua tingkat yang sama, pemisahannya $2|\beta|$ dan kedua elektron
$\pi$ memperoleh $2|\beta|$, yang sebanding dengan $\cos\theta$. Nilai itu
turun menjadi separuhnya pada $\theta = 60^\circ$ dan menjadi nol pada
$90^\circ$.
\end{solution}

\begin{solution}{exo:b2:fragment-orbitals:9}
Boran mempunyai enam elektron valensi: elektron-elektron itu mengisi
ketiga \omterm{def:b2:diatomic-mos:bonding}{orbital ikatan}, dan orbital p boron yang tegak lurus bidang tetap
kosong; tidak ada yang diperoleh dengan menjadi piramidal. Amonia
mempunyai dua elektron lebih banyak: keduanya menempati orbital yang
bersifat nonikatan pada bentuk datar, dan, seperti pada air, bentuk
piramidal memungkinkan orbital itu bercampur dengan orbital 2s dan
kombinasi hidrogen lalu turun.
\end{solution}

\begin{solution}{exo:b2:fragment-orbitals:10}
Persamaan sekular untuk $c_1 = c_2 = c_3$ memberikan $\alpha + 2\beta$; kedua
kombinasi yang ortogonal terhadapnya memberikan $\alpha - \beta$
(terdegenerasi). Dua elektron dalam $\alpha + 2\beta$: energinya
$2\alpha + 4\beta$, lebih rendah daripada $\ce{H2} + \ce{H+}$
($2\alpha + 2\beta$) karena $\beta < 0$.
\end{solution}

\begin{solution}{exo:b2:fragment-orbitals:11}
$\pi_1$: ketiga koefisien bertanda sama, tanpa simpul (ikatan); $\pi_2$:
tanda berlawanan pada karbon-karbon ujung dan sebuah simpul di tengah
(nonikatan); $\pi_3$: tanda berselang-seling, dua simpul (antiikatan).
Anionnya mempunyai empat elektron $\pi$: $\pi_2$ adalah tingkat terisi
tertingginya, dengan kerapatan hanya pada karbon-karbon ujung.
\end{solution}

\begin{solution}{exo:b2:fragment-orbitals:12}
Pada setiap dari kedua bidang yang memuat sumbu, tiga orbital p yang
sejajar (O, C, O) menghasilkan satu \omterm{def:b2:diatomic-mos:bonding}{orbital ikatan}, satu \omterm{def:b2:diatomic-mos:bonding}{orbital nonikatan}
(hanya pada oksigen), dan satu \omterm{def:b2:diatomic-mos:bonding}{orbital antiikatan}: seluruhnya enam orbital
$\pi$. Kedua \omterm{def:b2:diatomic-mos:bonding}{orbital ikatan} dan kedua \omterm{def:b2:diatomic-mos:bonding}{orbital nonikatan} terisi: delapan
elektron $\pi$ (dua ikatan $\pi$ dan dua pasangan elektron bebas oksigen
dalam gambaran Lewis).
\end{solution}

\begin{solution}{pb:b2:fragment-orbitals:1}
\textbf{1.} Delapan: enam dari oksigen, satu dari setiap hidrogen.
\textbf{2.} $h_1 + h_2$ dan $h_1 - h_2$.
\textbf{3.} 2s dan $2\mathrm p_z$.
\textbf{4.} $2\mathrm p_y$.
\textbf{5.} $2\mathrm p_x$, yang tegak lurus bidang molekul.
\textbf{6.} Enam; empat terisi.
\textbf{7.} $h_1 - h_2$, yang berganti tanda, seperti $2\mathrm p_z$, melalui
bidang yang tegak lurus sumbu pada oksigen.
\textbf{8.} $h_1 + h_2$.
\textbf{9.} $2\mathrm p_x$ dan $2\mathrm p_y$, yang tegak lurus sumbu: sepasang
\omterm{def:b2:diatomic-mos:bonding}{orbital nonikatan} yang terdegenerasi.
\textbf{10.} (2s, $h_1 + h_2$) ikatan$^2$, ($2\mathrm p_z$, $h_1 - h_2$) ikatan$^2$,
lalu $2\mathrm p_x^2\,2\mathrm p_y^2$.
\textbf{11.} Tidak: pasangan yang terdegenerasi itu penuh, setiap elektron
berpasangan.
\textbf{12.} Pada \omterm{def:b2:atomic-orbitals:orbital-energy}{energi orbital} 2p oksigen: tingkat itu nonikatan.
\textbf{13.} $2\mathrm p_y$ (dengan sumbu molekul linear sepanjang $z$ dan
pembengkokan pada $yz$).
\textbf{14.} Orbital itu bercampur dengan 2s dan $h_1 + h_2$ lalu turun:
tingkat $3a_1$.
\textbf{15.} Tingkat ikatan ($2\mathrm p_z$, $h_1 - h_2$), yang kehilangan
sebagian tumpang-tindih ketika hidrogen-hidrogen meninggalkan sumbu:
tingkat itu menjadi $1b_2$.
\textbf{16.} Tingkat terisi yang turun memperoleh lebih banyak daripada yang
hilang dari tingkat yang lain: dengan delapan elektron, bentuk bengkok
lebih stabil.
\textbf{17.} $104.48^\circ$, di antara $90^\circ$ (ikatan yang hanya dibuat dari
orbital p) dan $109.5^\circ$: orbital 2s ikut berperan dalam ikatan.
\textbf{18.} $2\mathrm p_x$, yang tetap menjadi $1b_1$ nonikatan.
\textbf{19.} Empat.
\textbf{20.} Dari $1b_1$, orbital $2\mathrm p_x$ oksigen yang nonikatan.
\textbf{21.} Mengambil elektron nonikatan hampir tidak mengubah
ikatan-ikatannya: ion itu mempunyai geometri molekul semula dan hanya
sedikit energi vibrasi yang tereksitasi. Mengambil elektron ikatan
mengubah panjang dan sudut ikatan, dan pitanya melebar meliputi banyak
tingkat vibrasi.
\textbf{22.} \qty{12.62}{eV} dibandingkan dengan \qty{13.62}{eV}: oksigen
dalam air membawa muatan parsial negatif yang ditarik dari
hidrogen-hidrogen, yang menaikkan tingkat-tingkat 2p-nya.
\textbf{23.} Keduanya bukan dua tingkat yang sama: kedua kombinasinya
adalah orbital $3a_1$ dan $1b_1$, yang energinya berbeda. Kedua gambaran
memerikan kerapatan elektron total yang sama; hanya orbital-orbital
terdelokalisasi yang memberikan energi ionisasinya.
\textbf{24.} Air mempunyai \textbf{empat} pita ionisasi valensi, yang
terendah pada $\boldsymbol{\approx \qty{12.6}{eV}}$.
\end{solution}
